% 边缘相同而联合分布不同的两种配对。教学数据：Z 等概率取 -1 与 1。
% (a) (Z,Z)：联合质量在对角 (-1,-1),(1,1) 各 1/2，坐标和 2Z 取 -2,2 各 1/2；
% (b) (Z,-Z)：联合质量在反对角 (-1,1),(1,-1) 各 1/2，坐标和恒为 0；
% 两种配对的行边缘与列边缘都是 (1/2,1/2)。
\begin{tikzpicture}[line width=.6pt,font=\fontsize{8}{10}\selectfont,>=Stealth,
  cell/.style={draw=FigureStroke,minimum width=11mm,minimum height=6.5mm,inner sep=0pt},
  marg/.style={cell,fill=FigureGreenFill},
  hl/.style={cell,fill=FigureBlueFill},
  hr/.style={cell,fill=FigureRedFill}]
\foreach \pane/\shift/\tag/\a/\b/\c/\d/\st/\sumtext/\ttl in {
   a/0/{(Z,Z)}/0.5/0/0/0.5/hl/{和为$2Z$：取$\pm2$各$1/2$}/{(a) 同号配对},
   b/60/{(Z,-Z)}/0/0.5/0.5/0/hr/{和恒为$0$}/{(b) 反号配对}}{
 \begin{scope}[xshift=\shift mm]
  \node[FigureInk] at (11mm,13.5mm) {$\tag$};
  \foreach \j/\y in {0/-1,1/1}{\node[FigureInk] at (11mm*\j+5.5mm,7mm) {$y=\y$};}
  \node[FigureGreen] at (30mm,7mm) {边缘};
  \node[FigureInk,anchor=east] at (-1.5mm,0) {$x=-1$};
  \node[FigureInk,anchor=east] at (-1.5mm,-6.5mm) {$x=1$};
  \ifdim\a pt>0pt \node[\st] at (5.5mm,0) {$\a$};\else\node[cell] at (5.5mm,0) {$0$};\fi
  \ifdim\b pt>0pt \node[\st] at (16.5mm,0) {$\b$};\else\node[cell] at (16.5mm,0) {$0$};\fi
  \ifdim\c pt>0pt \node[\st] at (5.5mm,-6.5mm) {$\c$};\else\node[cell] at (5.5mm,-6.5mm) {$0$};\fi
  \ifdim\d pt>0pt \node[\st] at (16.5mm,-6.5mm) {$\d$};\else\node[cell] at (16.5mm,-6.5mm) {$0$};\fi
  \node[marg] at (30mm,0) {$0.5$};
  \node[marg] at (30mm,-6.5mm) {$0.5$};
  \node[marg] at (5.5mm,-16mm) {$0.5$};
  \node[marg] at (16.5mm,-16mm) {$0.5$};
  \node[FigureGreen,anchor=east] at (-1.5mm,-16mm) {边缘};
  \node[FigureMuted,font=\fontsize{7.5}{9}\selectfont] at (14mm,-24mm) {\sumtext};
  \node[FigureInk] at (14mm,-31mm) {\ttl};
 \end{scope}}
\end{tikzpicture}
